\section{Introduction}
%context
During the past two years, the field of artificial intelligence (AI) has witnessed an explosion in public interest and economic investment. This has brought incredible breakthroughs and a constant flow of innovation. Its novel and disruptive nature, rate of improvement, and general applicability to every branch of science, education, and business make this technology unavoidable, and a good subject of study in the field of business engineering. Although a hardly understood field, it has become clear that artificial intelligence is set to change human experience considerably. \\
Data scientists strive to bridge the gap between research and productive environments by adapting this technology to all use cases. What was first commonly understood as chatbots or conversational agents able to generate instant responses to queries, large language models (LLM) are now being expanded to more complex systems that show a certain degree of agency and retrospection. This becomes useful for tackling multi-step unstructured problems with unpredictable outcomes. This area of artificial intelligence englobes paradigms and concepts built on top of LM technology, called agentic technology. The outcome of this technology is generative agents. \\ %goal
Researchers regularly struggle to keep up with newly published literature on artificial intelligence in general, and generative agents are not an exception. For this reason, this seminar paper aims to synthesize the findings regarding generative agents and assess the potential implications for the field of business engineering. The research is based on the following investigative question: \\
    
    \textit{What application can generative agents have inside productive organizations?} \\\\ %procedure
This paper does not entertain the various philosophical discussions on intelligence and consciousness, nor the ethical or moral implications that come with the dissemination of this technology. Rather, generative agents are looked at through a business engineering perspective by describing the fundamentals of this technology and assessing its economic and business impact.\\
Starting with the broader definition of artificial intelligence to later focus on the definitions and concepts behind agentic technology, this work reviews the most relevant and actual publications on generative agents. Once defined, the state-of-the-art is analyzed and capabilities and limitations are extracted. \\
It is widely assumed that the nature of business intelligence will change with artificial intelligence, but there is a lack of understanding of how and to what extent. The hope is to shine a light on the gap between the present and the future.
\section{Review of the Research in Generative Agents}
Although the literature on this area of technology has a short lifespan, it is nevertheless important to periodically review new findings to have an actualized overview of the field and the direction progress takes. In computer science, natural selection is ruthless towards new paradigms. Only the best ideas are implemented, transforming new technology into old. This paper makes the case for generative agents to be the safer bet as the next step towards artificial general intelligence (AGI). \\
In his book \textit{Thinking, fast and slow}, Nobel prize laureate Daniel \textcite{kahneman_thinking_2012} distinguishes two systems inside the human brain, \textit{system 1} and \textit{system 2}. System 1 reacts to specific cues with interiorized reactions, e.g., breathing, or with pre-trained behavior, e.g., fastening the seatbelt after entering a vehicle. System 2 on the other hand, represents rational behaviors, such as the ability to discern, reflect, or make more complex decisions. The brain is to be understood as a complex system with multiple simple specialized systems for single tasks, and one general cortex able to tackle problems in a more general sense. \\
Artificial intelligence has achieved historic milestones, e.g. surpassing the Turing test, cracking protein-folding algorithms, or beating world-class game players, by emulating the narrow and specific system 1 nature. This chapter leaps through disruptions leading to generative agents, a more system-2-type approach to tackle problems of a more general nature.
\subsection{Fundamentals of Artificial Intelligence}
Defining artificial intelligence becomes challenging, as its meaning has been diluted by its abuse in public discussion --mainly for marketing purposes. In broader terms, AI means the exhibition of intelligent behavior by a machine or a computer. It encompasses many research areas including machine learning, content generation, pattern recognition, and natural language processing \parencite[cf.][pp. 1-2]{russell_artificial_2010}. \\\\
\textbf{Generative AI}\indent Machine learning models can be divided into predictive and generative. Predictive machine learning models are the ones able to infer or extract new insights from a vast and unstructured dataset. Generative AI, on the other hand, is a model that has been trained on a carefully curated and structured dataset, so it can later imitate the information inside and create (\textit{generate}) new information \parencite[cf.][p. 112]{feuerriegel_generative_2024}.\\\\
\textbf{Large Language Models}\indent Language Models are a subset of generative AI capable of processing natural language. This kind of model is trained to predict the most probable next word (\textit{token}) given an initial input (\textit{prompt}). This is achieved not only by learning knowledge and facts about the world but also by the syntactical structure of human language. However, these statistical models can only store probabilities and not complete clear facts, leading to a variety of problems.\\
The ability to redact text and store facts leads users to anthropomorphize these models when in reality they present rather limited reasoning and logical skills. It has been demonstrated that by scaling the training data, higher reasoning abilities emerge \parencite{zeng_skywork-math_2024}. Language models trained on large datasets with a large number of parameters receive the name of \textit{large language models (LLMs)}. When an LLM is surrounded by a framework of applications that equip the model with additional capabilities outside language processing, they receive the name of \textit{foundational models} \parencite[cf.][p. 82]{Amaratunga.2023}. \\
Scaling, however, comes with an enormous cost, because training a large language model requires amounts of electricity and a server structure only available to the most important software companies in the world. Moreover, logical reasoning skills, especially mathematical abilities, seem to be immune to scaling. \\\\
\textbf{Transformer Architecture} \indent Researchers must, therefore, discover new designs for large language models. Transformer architecture revolutionized the way language models can infer meaning from words regardless of their position in the sentence, something especially useful for translation and human-like text generation \parencite{Vaswani.12.06.2017}. \\\\
\textbf{Prompt Engineering}\indent By spending time with these systems, experts have also discovered that by reformulating the prompt, the model can not only show more desired responses but also increase its reasoning abilities. This technique is called \textit{prompt engineering}. Most notable is the addition of the sentence \textit{"let's think step by step"} at the end of a complex prompt, e.g. math problem. Somehow the model deconstructs the problem into a series of smaller problems and follows a logical path to a solution \parencite[cf.][pp. 3-4]{Kojima.24.05.2022}. \\
This approach to prompt engineering is known as \textit{zero-shot chain-of-thought prompting}. Zero-shot means that no previous examples of a correct answer are given to the model, while chain-of-thought refers to triggering the model to think step by step. Prompt engineering has proven to be key in agentic technology, for the abilities of language and reason are deeply intertwined.\\\\
\textbf{Retrieval Augmented Generation (RAG)}\indent  The limit on reasoning capabilities is not the only obstacle large language models face. Another one is the inability to learn new information and access proprietary data. The model will produce a syntactically correct but factually fabricated answer when asked to infer knowledge not contained in its parameters. This is called a \textit{hallucination}.\\
As previously mentioned, training a model is expensive, and is only done once. It can later be fine-tuned to adapt it to certain tasks or requirements, but is not something that can be done regularly. This means that the knowledge the model can remember is frozen in time from the moment training takes place. Moreover, LLMs are trained on freely available data on the internet which is sufficient for generating a convincingly human-like conversation, but unfit for most use cases that require business data.\\
Retrieval augmented generation is a framework that combines three activities, namely \textit{indexing}, \textit{retrieval}, and \textit{generation} to solve both hallucinations and access to proprietary data. The desired data is curated and divided into chunks and \textit{indexed} into a vector database. The user prompts are then \textit{indexed} into vector representations and search for similarity within the vector database. The content with the highest similarity will be added to the prompt, \textit{augmenting} its context. This augmented prompt will then be fed again to the LLM which \textit{generates} the final answer. This way the vector database can be inexpensively updated every time new data is required \parencite[cf.][p. 3]{huang_survey_2024}. \\\\
\textbf{Function Calling}\indent In case the user asks for information unknown to the model and also outside of the private data, the system should be able to perform a web search. In principle, LLMs can only produce text as output. However, special LLMs like \textit{toolformer} \parencite{schick_toolformer_2023} or \textit{granite} \parencite{abdelaziz_granite-function_2024} are precisely trained to detect when they need to perform certain actions and generate an API function call in JSON-format describing the action they want to be performed. This function is then interpreted by an external software that executes it and gives back the answer to the model. \\
Function calling goes beyond web searches. It can be applied to any other tools that perform actions language models can't, e.g. calculator, code generation, image generation, etc. The ability to use tools is a key enabler to agents, providing a sense of agency to perform actions and communicate with the environment and other agents. The development during the past three years of the technology described in this section has enabled the emergence of generative agents. The following section defines how this technology assembles into an agentic framework.
\subsection{Fundamentals of Agentic Technology}
Once the technological foundations on which generative agents stand are established, this section provides clarification and the latest in agentic technology. Although researchers have not yet agreed on a formal definition, \textcite[pp. 8-9]{jennings_roadmap_1998} define an \textit{agent} as a system that exhibits \textit{situatedness}, \textit{autonomy}, and \textit{flexibility}:
\begin{itemize}
    \item \textit{Situatedness}: To be constrained to, and interact with an environment by taking and processing inputs and expressing outputs. Such an environment can be spatial, e.g., for robotic agents, or virtual, e.g., the internet. The opposite of situatedness is disembodied intelligence, e.g., expert systems, which require a human to input the data and interpret the results.
    \item \textit{Autonomy}: An agent can \textit{autonomously} perform actions without human intervention, e.g., software daemons. Autonomous systems are also able to process feedback and learn from past experiences.
    \item \textit{Flexibility}: The ability to \textit{respond} and adapt to changes, \textit{proactively} act towards a goal without the need for an initial event (e.g., environmental changes or human intervention), and \textit{socially interact} with other artificial agents or humans, if required.
    \end{itemize}
This definition is correct and can be used to provide a general context. However, back then, they considered agents a broader kind of software system able to automate human processes, e.g., air traffic control systems, or even a thermostat. Nowadays, a more nuanced definition is required. \\\\
\textbf{LLM-based Agents} \indent Research on agents is hardly a new field but is undergoing a resurgence due to recent advancements in LLM technology. Using language models to build agents stems from the \textit{"ReAct"} paradigm introduced by \textcite{Yao.06.10.2022} (the word is originated by combining the verbs \textit{reflect} and \textit{act}). In essence, this mechanism allows LLM chatbots to have a degree of introspection, by maintaining a conversation with themselves,  and deciding whether to use tools or perform other actions. This dichotomy is useful for automating multi-step tasks that require constant analysis and reaction to feedback.\\
Building on top of \textit{ReAct}, \textcite[pp. 8-11]{Park.07.04.2023} provide a framework for behavior in an open world, where the agent must perceive its environment and react accordingly. Their agents are equipped with a memory stream, where past experiences or previously retrieved information can be recalled into context and taken into account for further actions. In addition to \textit{reflect} and \textit{act}, these can also \textit{perceive} and \textit{plan}, and their output is considered to be a \textit{behavior} (see Figure~\ref{fig:agent}). The memory stream follows an algorithm to prioritize memories according to \textit{recency} and \textit{importance}. \\ \noindent
\begin{figure}
    \centering
    %\includegraphics[width= 1 \textwidth]{GenAgents.jpg}
    \caption{Generative Agent Framework, \textcite{Park.07.04.2023}}
    \label{fig:agent}
\end{figure}
Periodically, higher-level abstractions are generated reflecting on the output of the main memory lane and stored in a second memory, which allows for goal orientation. Otherwise, agents often lose track of the original objective when prompted with complex tasks. Additionally, the ability to plan becomes crucial to maintain the agent's behavior consistent over time. When planning, the system divides the task by describing a sequence of events that need to be accomplished. This plan is available for retrieval throughout the agentic process. \\\\
\textbf{Multi-Agent Architecture} \indent The natural step forward is the horizontal scaling of agentic technology by making different agents collaborate towards a common goal. A formal definition provided by \textcite[pp. 8-9]{wolpert_introduction_1999}, states that a multi-agent system must involve:
\begin{enumerate}
    \item The \textit{decomposition} of a complex task into multiple simpler subtasks, and each one is assigned to different agents.
    \item A \textit{communication} channel is established between agents to enable collaboration.
    \item A \textit{coordination} effort to integrate the separate efforts towards the common goal.
\end{enumerate}
Expanding this classical definition for today's agents, agent profiling is the key concept allowing for \textit{decomposition}. Each agent is assigned a role that possesses a global and a particular component. For example, for a software development multi-agent system, the global component defines the system's main objective (in this case, software development). The particular component speaks to the specific role performed by a single agent, e.g., developer, or scrum master \parencite[cf.][p. 3]{guo_large_2024}.\\\\
Currently, the research community focuses its efforts on solving the \textit{communication} problem. Classical communication protocols are faced with a trade-off between the versatility of the channels to support different media formats, the computational load of such system, and the portability (or adaptability) to different environments and agent types \parencite[p. 3]{marro_scalable_2024}. \\
Papers proposing innovative communication algorithms are being published regularly. One example is \textit{Agora}, a meta-protocol published by \textcite[pp. 5-6]{marro_scalable_2024} that employs already existing protocols in a hierarchical manner. This system can adapt according to the circumstances. For instance, a certain protocol might be preferred if a situation requires high versatility. \\
\textit{Agora} perfectly illustrates how this technology dynamically builds on top of previous technology, exponentially scaling capabilities. The authors remark on how little human intervention is required at the implementation --a mere indication to the agents which tools they have access to sufices.\\
Regarding \textit{coordination}, \textcite[pp. 16-17]{li_metaagents_2023} have found that the reasoning and reflecting modules, when applied to retrieved conversations with other agents, enhance the collaborative effort. However, they detected a \textit{misplacement} problem when agents assign tasks within teams. This is caused by an exaggeration of the agents when describing their capabilities or skills, causing suboptimal results. This problem is currently being investigated, but seems to be a reminiscence of the \textit{hallucinating} issues primitive large language models faced. For this reason, multi-agent systems only exist in research scenarios and are not deployed in production.
\section{Generative Agents in Business Enginering}
Regardless of the limitations, the previous chapter outlined how researchers are able to organize teams of agents that can solve not only tasks but complex projects like developing software. Currently, this affects mostly knowledge tasks but one could envision a future in which agents acquire motor skills and can have a similar kind of cooperation in a spatial setting. Hardware developments are underway. Small, portable supercomputers able to host a small LLM locally are already commercially available. In the foreseeable future, every machine could have attached one of these \textit{"reasoning modules"}, hosting an LLM-based agent fully integrated with its motor capabilities. \\
Regardless, the amount of work that could potentially be done exclusively in the intellectual realm is outstanding. In the next section, the capabilities and limitations are further discussed and contextualized into a business setting.
\subsection{Opportunities and Capabilities}
\textcite[pp. 5-10]{guo_large_2024} distinguish two main classes of LLM-based multi-agent systems capabilities with their subsequent possible applications:
\begin{itemize}
    \item \textit{Problem-Solving}: Software development, embodied agents, science experiments, psychological experiments, etc.
    \item \textit{World Simulation}: Societal simulation, psychology, game theory, economic and financial trading simulations, recommender systems, policy making, propagation simulation, among others.
\end{itemize}
A concrete example of an application for this technology is the simulation of human behavior in a virtual setting. Companies can gather user information and recreate their habits into a virtual persona to better understand their needs. This could be then used to execute tailor-made services that facilitate day-to-day experiences \parencite[cf.][p. 17]{Park.07.04.2023}. \\
Brought to a business context, a virtual assistant can trigger customized alarm clocks, the automatic execution of devices or programs, and the management of events and tasks, to name a few. This could be understood as having a clone in the virtual realm (or a \textit{digital twin}) that represents and acts on the worker's behalf.\\\\
On the business side, it is hard to find literature on the impact of generative agents, due to their young nature. Nonetheless, GenAI's economic impact has been more broadly addressed. In a recent study published in the \textit{Review of Managerial Science}, \textcite[pp. 1211-1213]{kanbach_genai_2024} have studied and divided GenAI's impact across industries into three distinctive categories:
\begin{itemize}
    \item \textit{Value Creation Innovation}: GenAI facilitates new capabilities, technologies, partnerships, and processes by reshaping information access, content creation, and business operations. It offers efficiency improvements and the potential for completely novel products and services.
    \item \textit{New Proposition Innovation}: GenAI can lead to new offerings, markets, channels, and customer relationships, particularly affecting white-collar knowledge workers and shifting human roles from creators to editors.
    \item \textit{Value Capture Innovation}: GenAI can drive new revenue models and cost structures, reducing content production costs and enabling mass customization.
\end{itemize}
The authors emphasize that GenAI will continue to evolve and permeate various industries, pushing the boundaries of what is possible. Businesses must adapt to harness GenAI's power, with early adopters potentially gaining a competitive advantage. However, significant challenges and ethical concerns, such as biases, disinformation, and intellectual property theft, must be addressed through ethical frameworks, transparency, and new policies. \\\\
In the case of Germany, a report published by \textcite{mckinsey2023genai} explores how GenAI can address the pressing issue of skilled labor shortages while driving productivity and economic growth with skilled labor shortages affecting approximately 50\% of businesses by 2022 (a fivefold increase since 2009). Open job positions have quadrupled during this period, emphasizing the need for innovative solutions.\\
McKinsey suggests that GenAI has the potential to mitigate these shortages by automating tasks and enhancing productivity across various sectors. Unlike traditional analytical AI, which focuses on classifying and predicting data, GenAI generates creative outputs, significantly augmenting human capabilities. GenAI has become especially suitable for fields that traditionally were out of automation's reach, such as education, STEM, or healthcare.\\
The report highlights that GenAI's automation potential is most pronounced for individuals with tertiary education, but its societal benefits also extend to high school graduates in fields like community health and technical support. Moreover, the highest wage earners stand to benefit most, as these roles often involve complex, non-repetitive tasks well-suited for \textit{augmentatiion} by GenAI. \\
Conclusively, experts at McKinsey deliver recommendations to fully leverage GenAi. These \textit{key enablers} are:
\begin{itemize}
    \item \textit{Skill Development}: Upskilling and reskilling the workforce to adapt to new roles created by AI technologies.
    \item \textit{Supportive Ecosystem}: Establishing the right infrastructure, policies, and investment frameworks to foster AI innovation and adoption.
\end{itemize}
Germany has strong AI skill foundations, ranking second among OECD countries, and its significant number of GenAI startups positions it as a leader in this field. However, the report also notes that funding and investment levels in AI projects need to increase to maintain this competitive edge.
\subsection{Risks and Limitations}
In the previous section, the technological limitation of \textit{misplacement} was discussed. Due to the complexity and early lifespan of these systems, this is hardly the only one. On the business side, challenges and ethical concerns such as biases, disinformation, and intellectual property theft are also points of concern. \\
By reviewing \textcite[p. 18]{li_metaagents_2023} and \textcite[pp. 10-11]{guo_large_2024}, the following main technological limitations are distinguished:
\begin{itemize}
    \item \textit{Scaling Up}: Scaling the number of agents in a system currently requires vast amounts of computation, memory, and power. Also, scaling the timespan of the simulations is necessary for any real-life project. Furthermore, adding agents to a system brings orchestration issues that still lack a feasible solution.
    \item \textit{Elevating Complexity}: These systems exist in very controlled scenarios, where the amount of possible inputs and variability is seriously reduced compared to base reality.
    \item \textit{Advancing into Multi-Modal Environments}: Integrating diverse data types requires more complex communication protocols. As discussed previously, these require considerable hardware specifications and energy consumption.
    \item \textit{Addressing Hallucinations}: This issue speaks to the aforementioned \textit{misplacement} issue, where models cannot distinguish factual reality from fiction, and are poorly trained to acknowledge ignorance or self-deficiencies. 
    \item \textit{Acquiring Collective Intelligence}: Agents are currently adjusted in isolation, with communication protocols as the only integrative structure within the system. Simultaneously adjusting multiple agents still poses a challenge for experts.
\end{itemize}
Progress in data and infrastructure availability is also missing. Before considering this technology as a widespread economic factor, the necessary infrastructure must be set in place. For instance, the development of 5G networks, the international cooperation for building extensive training datasets, the access to more energy sources, and competent data centers that can handle this technology on a massive scale. \\
Also, a regulatory framework must be drafted and implemented to balance AI's transformative potential with ethical considerations and worker's rights, emphasizing a \textit{"human-centered" AI} that focuses on the well-being, dignity, and preserving the right for self-determination of citizens. Recognizing AI's potential to enhance competitiveness while remaining vigilant about its associated risks, e.g. job displacement, privacy, and surveillance concerns, etc. \parencite[cf.][p. 59-66]{krzywdzinski_promoting_2023}. \\
However, there is a risk of overregulation, which might help mitigate the short-term impact of AI but can decrease economic growth in the mid-to-long term, which would negatively affect labor. One example of this dichotomy has been the strikes in American ports demanding higher wages and a restriction to automation. It is hard to assess the implications of complying with unions and purposely neglecting innovation, but the dangers of having obsolete economic sectors (in this case ports) should not be understated.
\section{Conclusion}
After reviewing the literature, it becomes evident that experts lack a definitive perspective on the impact of agentic technology in a business setting. LLM-based multi-agent systems represent a significant leap in the automation of intellectual tasks. These systems can emulate fundamental human cognitive capacities, such as reasoning, problem-solving, and generating creative outputs, making them invaluable for tasks like data analysis, content creation, and strategic planning. Additionally, their ability to communicate and collaborate enables them to function as cohesive units that can address complex, multifaceted challenges. \\
However, their performance is bounded by the quality of training data, limitations in interpretability, and the inability to handle tasks requiring fine-grained motor skills or physical-world interactions, which remain the domain of human expertise. The pace of improvement in these systems suggests that current constraints might diminish rapidly, but the challenge of generalizing their applications beyond intellectual tasks persists. On the other hand, LLM-based multi-agent systems are constrained by significant practical barriers, including the need for vast computational resources and energy consumption, which pose scalability challenges. \\
Furthermore, while their potential for learning and adaptation is impressive, they still lack the innate human capacity for ethical reasoning, emotional intelligence, and contextual awareness in real-world scenarios. This gap emphasizes the necessity of human oversight and hybrid systems where AI supports, rather than replaces, human decision-making.\\
The concern for the potential loss of the intrinsic meaning derived from professional roles remains. Some argue that the era of craftsmanship—a key pillar of Germany's industrial identity—is fading for many professions. \\
A more positive framing of this problem is to think about how much more humanity will be able to achieve if all kinds of mindless and repetitive tasks can be delegated to computers, enabling workers to concentrate on innovation, creativity, and skilled labor. This could usher in an era of prosperity, where technological advancements empower existing workers and mitigate the lack of high-skilled professionals. In the future, humans might find meaning in social contributions or artistic endeavors, even if there is no financial incentive. This behavior has already been shown throughout history by individuals who possess extreme abundance.\\
Ultimately, even if all the previous assumptions are true and the future becomes fully automated, it is the job of the business engineer to look at the other side of the coin and always view the glass half full. Behind every system there is an administrator, and the more complex the system, the more maintenance is required. Moreover, the future has to be built first, and business engineering will be the stepping stone that bridges the present with the future.
